\chapter{Метод, модель и программная реализация}

\section{Математическая модель квадрокоптера}

За основу взята нелинейная модель квадрокоптера из диссертации~\cite{kim_algorithms}, после чего она была адаптирована под численное моделирование и под связку с регулятором траекторного слежения. Базовая кинематика записывается в виде
\begin{gather}
    \dot{p} = v,
\end{gather}
\begin{gather}
    \dot{v} = \frac{b(\phi,\theta,\psi)\bigl(u_1 + g\bigr) - [0,0,g]^{\top}}{m},
\end{gather}
а вращательное движение задается через углы ориентации и соответствующие управляющие воздействия.

Для учета ограничений по тяге используется сглаженная функция насыщения
\begin{gather}
    u_1 = L \tanh\left(\frac{\bar{u}_1}{L}\right).
\end{gather}
Кроме того, в модель введены интеграторные расширения
\begin{gather}
    \dot{\bar{u}}_1 = \rho_1, \qquad \dot{\rho}_1 = v_1,
\end{gather}
\begin{gather}
    \dot{u}_2 = \rho_2, \qquad \dot{\rho}_2 = v_2,
\end{gather}
что позволяет привести систему к форме, удобной для реализации выбранного закона управления.

Практически это означает следующее: модель остается физически интерпретируемой, но при этом становится пригодной для прямой программной реализации в контуре управления и для серийной прогонки большого числа сценариев.

\section{Геометрия траектории и наблюдатель ближайшей точки}

Чтобы перейти от обычного описания движения к задаче слежения за кривой, для каждой точки полета нужно знать, какой точке опорной траектории она соответствует. Для этого используется наблюдатель ближайшей точки. Его работа основана на условии ортогональности ошибки положению касательного вектора:
\begin{gather}
    H(\zeta,p) = \bigl(p^{\ast}(\zeta) - p\bigr)^{\top} t(\zeta) = 0.
\end{gather}

Для прямых в проекте используется аналитическое вычисление ближайшей точки, а для окружностей и спиралей применяется итерационный наблюдатель. На практике именно этот блок связывает реальное положение объекта с параметром траектории и позволяет корректно вычислять ошибки в системе Френе.

По найденному параметру формируются геометрические характеристики кривой, которые далее используются и в регуляторе, и в нейросетевом модуле: направление касательной, локальная кривизна, угол между вектором скорости и касательной, а также длина уже пройденной дуги.

\section{Реализация базового регулятора}

Регулятор траекторного слежения строится по регулируемому выходу $\lambda_1$ и использует расширенный наблюдатель производных высокого усиления. В программной реализации это позволяет обойтись без явного численного дифференцирования шумных сигналов и сохранить структуру алгоритма, предложенную в исходной теоретической работе.

Общая идея закона управления состоит в том, что сначала по текущим ошибкам и их оценкам формируется виртуальный управляющий сигнал, а затем он преобразуется в реальные управляющие воздействия объекта. Для этого используются матрицы преобразования, учитывающие одновременно геометрию траектории и ориентацию квадрокоптера.

Отдельно была доработана первая компонента ошибки синхронизации. Если использовать выражение $V^{\ast} t$, то при адаптивном изменении скорости в системе возникает скачок ошибки, который плохо сочетается с наблюдателем высокого усиления. Поэтому в реализации используется интегральное накопление опорной дуги:
\begin{gather}
    s_{\mathrm{ref}}(t) = \int\limits_0^t V^{\ast}(\tau)\, d\tau.
\end{gather}
Такое решение оказалось принципиально важным: оно сохраняет непрерывность опорного сигнала и позволяет безопасно подключать адаптивный выбор скорости.

\section{Проверка корректности базового контура}

Перед подключением интеллектуального модуля был проведен отдельный этап проверки базового регулятора. Для этого подготовлен набор автоматизированных сценариев на разных типах кривых: прямая, окружность, спираль и несколько дополнительных конфигураций с другой кривизной и наклоном.

Проверка строилась по простому и понятному критерию: к концу моделирования вектор ошибок должен оставаться в допустимой окрестности нуля, а само моделирование не должно уходить в неустойчивый режим. Такой этап был нужен не формально, а по сути, потому что без устойчивого базового контура дальнейшее исследование адаптивного выбора скорости теряло бы смысл.

\section{Формирование датасета для выбора скорости}

Нейросетевой модуль в работе не управляет объектом напрямую. Его задача уже уже: по текущему состоянию и локальной геометрии траектории оценить максимально допустимое значение параметрической скорости $V^{\ast}$. Поэтому для обучения сначала формируется датасет примеров вида <<состояние - допустимая скорость>>.

Каждый обучающий пример строится следующим образом. Выбирается одна из допустимых кривых, затем на ней берется опорная точка, вокруг которой формируется возмущенное начальное состояние. После этого из состояния извлекается вектор признаков
\begin{gather}
    f =
    \begin{bmatrix}
        e_1 &
        e_2 &
        \dot{e}_2 &
        \|v\| &
        e_{\mathrm{head}} &
        \kappa_c &
        \kappa_{\max}
    \end{bmatrix}^{\top}.
\end{gather}

В этот набор вошли признаки, которые действительно влияют на допустимую скорость: текущее отклонение от траектории, скорость изменения ошибки, фактическая скорость движения и характеристики геометрии ближайшего участка кривой. Такой выбор получился достаточно компактным, но при этом информативным.

Целевое значение скорости определяется оракулом. Он перебирает кандидаты в диапазоне
\begin{gather}
    V \in [V_{\min}, V_{\max}]
\end{gather}
и выбирает наибольшее значение, при котором короткий rollout остается устойчивым:
\begin{gather}
    V_{\mathrm{opt}} = \max \left\{ V \mid \mathrm{Stable}(x_0,\mathrm{curve},V) = True \right\}.
\end{gather}
По сути, датасет фиксирует опыт базового контроллера: в каких локальных условиях система еще может двигаться быстрее, а в каких уже нужно снижать темп.

\section{Нейросетевой модуль}

Для аппроксимации зависимости между признаками и допустимой скоростью используется полносвязная нейронная сеть `SpeedMLP` со структурой
\begin{gather}
    7 \rightarrow 128 \rightarrow 128 \rightarrow 64 \rightarrow 1.
\end{gather}
На скрытых слоях применяется функция активации ReLU, а на выходе используется ограничение, не позволяющее модели предсказывать скорость вне допустимого диапазона.

Обучение проводится как задача регрессии по среднеквадратичной ошибке:
\begin{gather}
    \mathcal{L}_{\mathrm{MSE}} =
    \frac{1}{N}\sum_{i=1}^{N}\left(\hat{V}^{\ast}_i - V_{\mathrm{opt},i}\right)^2.
\end{gather}
Во время обучения используется разбиение на обучающую и валидационную выборки, а также ранняя остановка. Это позволяет не просто подогнать сеть под набор данных, а получить модель, которая сохраняет адекватное поведение на новых примерах из того же семейства траекторий.

\section{Интеграция нейросети в контур управления}

После обучения нейросеть подключается как внешний модуль, который на каждом шаге моделирования выдает новое значение целевой скорости. При этом она не получает полной свободы. В программной реализации добавлены ограничения, которые нужны для безопасной работы:
\begin{enumerate}
    \item скорость ограничивается сверху и снизу;
    \item ограничивается темп изменения $V^{\ast}$;
    \item в аварийных режимах нейросети запрещается увеличивать скорость;
    \item на начальном участке используется режим прогрева с постоянной базовой скоростью.
\end{enumerate}

Такое решение отражает общую идею всей работы: нейросетевой модуль не подменяет классический регулятор, а аккуратно дополняет его. Устойчивость по-прежнему обеспечивается аналитическим контуром, а интеллектуальная часть только подбирает более удачный режим движения.

\section{Дополнительные архитектуры: SAC, TD3, PPO}
\label{sec:rl_models}

Помимо базовой архитектуры MLP, в работе исследованы три альтернативные архитектуры, построенные на принципах обучения с подкреплением, но применяемые в режиме обучения по готовому датасету~--- так называемом offline-режиме.

Все три модели принимают тот же вектор из семи признаков и предсказывают скалярную $V^{\ast}$. Отличие состоит в функции потерь и архитектурных деталях, которые меняют характер регрессии.

\textbf{SAC (Soft Actor-Critic, SpeedSAC).} Стохастический актор параметризует гауссово распределение $(\mu, \sigma)$ над $V^{\ast}$. Во время обучения используются два Q-критика для оценки качества предсказанного значения и энтропийный бонус, который предотвращает вырождение политики в детерминированную точку. Функция потерь:
\begin{gather}
    \mathcal{L}_{\mathrm{SAC}} = -\log \pi_\theta(V_{\mathrm{opt}} \mid f)
    - \alpha_{\mathrm{ent}} \cdot \mathcal{H}[\pi_\theta(\cdot \mid f)],
\end{gather}
где $\mathcal{H}$~--- дифференциальная энтропия, $\alpha_{\mathrm{ent}}$~--- коэффициент её веса. При инференсе возвращается среднее $\mu$~--- детерминированный выход.

\textbf{TD3 (Twin Delayed DDPG, SpeedTD3).} Детерминированный актор $a(f) \to V^{\ast}$ обучается одновременно по двум критериям: поведенческому клонированию (BC~MSE, минимизация отклонения от оракульных меток) и Q-guided корректировке с задержкой обновления. Два целевых критика с Polyak-усреднением стабилизируют обучение. Функция потерь актора:
\begin{gather}
    \mathcal{L}_{\mathrm{TD3}} = \mathcal{L}_{\mathrm{BC}} - \lambda \cdot \min_j Q_j(f, a(f)),
\end{gather}
где $\lambda$~--- вес Q-компонента.

\textbf{PPO (Proximal Policy Optimization, SpeedPPO).} Гауссова политика с отдельной оценочной головой обучается в оффлайн режиме поведенческого клонирования с ограничением PPO-clip и энтропийным бонусом:
\begin{gather}
    \mathcal{L}_{\mathrm{PPO}} = -\mathbb{E}\left[\min\!\left(
        r\, \hat{A},\;
        \mathrm{clip}(r, 1{-}\varepsilon, 1{+}\varepsilon)\,\hat{A}
    \right)\right]
    + c_v \mathcal{L}_v - c_{\mathrm{ent}} \mathcal{H},
\end{gather}
где $r = \pi_\theta / \pi_{\mathrm{old}}$, $\hat{A}$~--- оценка преимущества, $c_v$~--- вес ошибки оценочной функции.

Единый реестр моделей (\texttt{registry.py}) предоставляет общий интерфейс для создания, сохранения и загрузки любой из четырех архитектур. Класс \texttt{SpeedPredictorAny} обеспечивает взаимозаменяемый инференс: метод \texttt{predict(f)} работает одинаково для всех типов и возвращает $V^{\ast} \in [V_{\min}, V_{\max}]$.

\section{Модуль автоматизированной оценки}

Для воспроизводимого сравнения всех архитектур реализован отдельный модуль оценки (\texttt{ml/evaluation/}), работающий независимо от кода обучения. Он содержит три компонента.

\textbf{Тестовый набор кривых} (\texttt{test\_suite.py}) фиксирует четыре сценария с полными параметрами симуляции: спираль~$r{=}3$, горизонтальная окружность~$r{=}3$, спираль~$r{=}2$ и пространственная прямая. Набор детерминирован и воспроизводим~--- любой запуск бенчмарка работает с одинаковыми начальными условиями.

\textbf{BenchmarkRunner} (\texttt{benchmark.py}) принимает словарь путей к чекпоинтам моделей и для каждого сценария последовательно запускает: режим baseline (константная $V^{\ast}$) и каждую загруженную модель. Результаты упакованы в структуру \texttt{ModelResult} с метриками $e_{1,\mathrm{RMS}}$, $e_{2,\mathrm{RMS}}$, $e_{2,\max}$, $\bar{v}$ и коэффициентом ускорения.

\textbf{Модуль визуализации} (\texttt{plots.py}) строит временны́е профили ошибок и скорости для всех моделей на одном графике, а также сводные столбчатые диаграммы. Дополнительно генерируется \LaTeX-таблица, которая может быть напрямую включена в отчет.

Запуск полного бенчмарка:
\begin{verbatim}
python code/scenarios/run_benchmark.py
python code/scenarios/run_benchmark.py --models mlp,sac --curves spiral_r3
\end{verbatim}

\section{Выводы по главе}

Во второй главе описана полная схема программной реализации. Она включает нелинейную модель квадрокоптера, геометрический блок работы с траекторией, базовый регулятор слежения, процедуру формирования датасета, четыре архитектуры нейросетевого модуля выбора скорости и модуль автоматизированной оценки. Ключевым результатом этого этапа стала работоспособная многоуровневая архитектура, в которой аналитический контур отвечает за устойчивость, а нейросетевые компоненты~--- за адаптивный выбор режима движения.
